\documentclass{article}
\usepackage{tcee-repaso}
\tcetema[La PAC y la política pesquera común]{3.B.38}{La política agrícola de la Unión Europea. Problemas económicos y procesos de reforma. La política pesquera común.}
\tcefecha{19/11/2026}
\begin{document}
\cabecerarepaso
\begin{columnas}

% ==========================================================================
\bloque[2,5]{Introducción}
\begin{itemize}
\item \textbf{Enganche}: «Aunque la mona se vista de seda, mona se queda».
  \begin{itemize}
  \item La PAC cambia de traje (precios, hectáreas, desacoplamiento, pagos verdes); tractores en Bruselas, París y Madrid (2024).
  \item El 25\,\% de los perceptores se lleva el 85\,\% de los pagos directos; en el mar, sin reglas, el barco grande vacía el caladero.
  \end{itemize}
\item \textbf{Relevancia}: 386.600~M€ en el MFP 2021-2027 (casi 1/3) \vertema{3.B.37}; la política más antigua.
  \begin{itemize}
  \item España: 2.º perceptor ($\approx$7.400~M€/año) y 1.ª flota \vertema{4.A.4}.
  \end{itemize}
\item \textbf{Contextualización}: \concepto{PAC}: mercados, renta, desarrollo rural; \concepto{PPC}: recursos, flota, mercados, acuerdos.
  \begin{itemize}
  \item Compartida (art.~4~TFUE), salvo la conservación de los recursos marinos, exclusiva (art.~3) \vertema{3.B.36}.
  \end{itemize}
\item \textbf{Problemática}: ¿por qué intervenir y en común? ¿Por qué tantas reformas? ¿Qué queda tras 2027? ¿Cómo evitar la sobrepesca?
\item \textbf{Estructura}: \hyperref[rep:I]{I} fundamentos; \hyperref[rep:II]{II} reformas y problemas; \hyperref[rep:III]{III} pesca.
\end{itemize}
\esquemacante{%
\apartadocante{Introducción}{2,5}
\apartadocante{I. La PAC: fundamentos y configuración}{8}
\subcante{1. Justificación, naturaleza y objetivos}{3}
\subcante{2. Configuración inicial (1957-1992)}{5}
\apartadocante{II. Problemas y reformas de la PAC}{10,5}
\subcante{1. Reformas 1984-2020}{4,5}
\subcante{2. La PAC 2023-2027 y su simplificación}{2,5}
\subcante{3. Problemas y retos}{3,5}
\apartadocante{III. La política pesquera común}{8}
\subcante{1. Justificación, naturaleza y objetivos}{3,5}
\subcante{2. La PPC de 1983}{1}
\subcante{3. Reformas de 2002 y 2013}{1,5}
\subcante{4. Situación actual}{2}
\apartadocante{Conclusión}{1}}

% ==========================================================================
\bloque[8]{I. La PAC: fundamentos y configuración inicial\label{rep:I}}
\begin{ideaclave}
Hay razones de eficiencia y de equidad para intervenir, y de economía política para intervenir de más. La PAC de 1962 eligió el instrumento más distorsionador, el precio: abasteció a Europa a costa de consumidores, contribuyentes y terceros países.
\end{ideaclave}

\sub{I.1. Justificación, naturaleza y objetivos\label{rep:I.1}}
\subsub{a) Justificación económica}
\begin{itemize}
\item \emph{Eficiencia} \vertema{3.A.23}:
  \begin{itemize}
  \item externalidades (nitratos, erosión, metano; paisaje);
  \item bienes públicos (biodiversidad, seguridad del abastecimiento);
  \item inestabilidad de precios: oferta retardada, demanda rígida (telaraña \vertema{3.A.16}), seguros incompletos.
  \end{itemize}
\item \emph{Equidad: el «problema agrario»}:
  \begin{itemize}
  \item ley de \autor{Engel}, ley de \autor{King}; rueda del progreso técnico (\autor{Cochrane}, 1958);
  \item baja movilidad de factores: en 2017, algo menos de la mitad de lo que ganaría en otro empleo;
  \item equidad territorial.
  \end{itemize}
\item \emph{Economía política}: \autor{Olson} (1965): grupo pequeño y concentrado frente a consumidores dispersos; la protección crece con el desarrollo.
\item \emph{¿Por qué común?} Mercado común con precios sostenidos $\Rightarrow$ precios, protección y financiación comunes; pacto franco-alemán.
\end{itemize}
\subsub{b) Marco jurídico y toma de decisiones}
\begin{itemize}
\item Arts.~38-44 TFUE; \textbf{objetivos del art.~39}:
  \begin{itemize}
  \item productividad; nivel de vida equitativo; estabilizar los mercados;
  \item seguridad de abastecimientos; precios razonables.
  \end{itemize}
\item Art.~40~OCM; art.~42~competencia.
\item Art.~43: PLO, pero el Consejo fija precios, exacciones, ayudas, limitaciones y posibilidades de pesca.
\item Gestión compartida; en España, el FEGA.
\end{itemize}

\sub{I.2. Configuración inicial (1957-1992)\label{rep:I.2}}
\subsub{a) Origen y principios}
\begin{itemize}
\item Stresa (1958, \autor{Mansholt}): precios y mercados + estructuras.
\item 1962: OCM y \concepto{FEOGA} (Garantía $\approx$90\,\%, Orientación).
\item Principios: \concepto{unidad de mercado}, \concepto{preferencia comunitaria}, \concepto{solidaridad financiera}.
\end{itemize}
\subsub{b) Política de precios y mercados y de estructuras}
\begin{itemize}
\item Precios \emph{indicativo}, de \emph{intervención} (mínimo, compras públicas) y de \emph{umbral}.
\item \emph{Exacciones variables} y \emph{restituciones} \vertema{3.B.7}.
\item Agrimoneda: tipos verdes y montantes compensatorios (1969-1993). Plan \autor{Mansholt} (1968).
\end{itemize}
\subsub{c) El modelo del sostenimiento de precios}
\begin{itemize}
\item \emph{Idea}: ¿qué cuesta un precio interior por encima del mundial?
\item \emph{Supuestos} (\autor{Baldwin}-\autor{Wyplosz}): país pequeño, $p^I>p^W$, intervención, exacción y restitución $p^I-p^W$.
\item \emph{Desarrollo}:
  \begin{itemize}
  \item años 60, importa: productores $+a$, consumidores $-(a+b+c+d)$, exacción $+c$;
  \item revolución verde, exporta: productores $+(a+b+c)$, consumidores $-(a+b)$, restituciones $-(b+c+d)$:
\funcionclave{\Delta W=-(b+d)<0}
  \item país grande: hunde $p^W$; daño a terceros; exporta la inestabilidad.
  \end{itemize}
\end{itemize}
\graficorepaso[0.8\linewidth]{sostenimiento-precios}{De la exacción variable a las restituciones: los mismos triángulos de pérdida, $b$ y $d$}
\begin{itemize}
\item \emph{Implicaciones}: transferencia cara (6~a productores cuesta 8), regresiva; excedentes; conflicto en la Ronda Uruguay \vertema{3.B.33}.
\item \emph{Extensiones}: la explotación grande se lleva casi toda la ganancia; cuotas (renta capitalizada); \emph{set-aside} (\emph{slippage}).
\item \emph{Evidencia}: PSE de la OCDE 38,5\,\% (1986-88) $\to$ 18,2\,\% (2025).
  \begin{itemize}
  \item Solo $\approx$25\,\% del sostenimiento llega a renta (OCDE, 2003); pérdida $\approx$1\,\% del PIB a finales de los 80.
  \end{itemize}
\item \emph{Valoración}: regla de \autor{Tinbergen}: un precio para 5~objetivos; las reformas separan instrumentos.
\end{itemize}
\graficorepaso[0.85\linewidth]{pse-ue}{Apoyo al productor (PSE, OCDE): casi el 40\,\% de los ingresos agrarios a mediados de los 80, 18\,\% en 2025}
\subsub{d) Valoración de la PAC (1962-1992)}
\begin{itemize}
\item Productividad y abastecimiento, sí; renta y estabilidad, en parte; precios razonables, no.
\item Gasto: 73\,\% del presupuesto (1985); 80/20 en las ayudas; excedentes; conflictos; intensificación.
\end{itemize}

% ==========================================================================
\bloque[10,5]{II. Problemas económicos y procesos de reforma\label{rep:II}}
\begin{ideaclave}
Del precio al pago y del pago a las condiciones: fin de los excedentes y del conflicto comercial. Quedan el reparto, la capitalización en la tierra, el relevo, el medio ambiente y la legitimidad; tras 2027, ¿un fondo único por Estado?
\end{ideaclave}

\sub{II.1. Desarrollo, resultados y reformas (1984-2020)\label{rep:II.1}}
\subsub{a) Los primeros frenos y la reforma \autor{MacSharry} (1992)}
\begin{itemize}
\item 1984: cuotas lecheras (hasta 2015); 1988: estabilizadores y línea directriz (74\,\% del PNB).
\item \textbf{\autor{MacSharry}} (1992): precios $-$30\,\% cereales, $-$15\,\% vacuno.
  \begin{itemize}
  \item Pagos por ha y cabeza (\concepto{caja azul}); retirada de tierras 15\,\%; agroambientales.
  \item Blair House y Ronda Uruguay.
  \end{itemize}
\end{itemize}
\graficorepaso[0.95\linewidth]{gasto-pac}{Gasto de la PAC 1980-2024: del apoyo al mercado a los pagos acoplados, desacoplados y el desarrollo rural; de 0,66~a 0,31\,\% del PIB}
\subsub{b) El modelo del desacoplamiento}
\begin{itemize}
\item \emph{Idea}: 2.º teorema del bienestar, transferencia de suma fija \vertema{3.A.22}.
\item \emph{Supuestos}: se baja $p^I$ a $p^W$ y se compensa con $s$ por unidad o con un pago $D$ independiente de la producción.
\item \emph{Desarrollo}:
  \begin{itemize}
  \item bajar el precio: $+(b+d)$; con $D=a+b+c$, nadie pierde;
  \item acoplada: sigue $q_1$ y queda $d$:
\funcionclave{\text{pérdida}=\tfrac12 s^2q'(p)>0=\text{desacoplada}}
  \item eficiencia de la transferencia: 6/7 frente a 1.
  \end{itemize}
\end{itemize}
\graficorepaso[0.9\linewidth]{desacoplamiento}{Acoplada: más producción y triángulo $d$; desacoplada: sin pérdida}
\begin{itemize}
\item \emph{Implicaciones}: orientación al mercado, \concepto{caja verde}, \autor{Tinbergen}.
\item \emph{Extensiones}: desacoplamiento imperfecto: riqueza (\autor{Hennessy}, 1998), crédito, expectativas, salida, tierra; acopladas hasta 13 + 2\,\%.
\item \emph{Evidencia}: poco efecto sobre la producción (OCDE, 2001); más desinversión y salida (\autor{Kazukauskas} \emph{et al.}, 2013).
\item \emph{Valoración}: ¿por qué pagar por tener tierra? Condicionalidad y redistribución.
\end{itemize}
\subsub{c) Agenda 2000, \autor{Fischler} (2003) y chequeo médico (2008)}
\begin{itemize}
\item Agenda 2000: \concepto{segundo pilar} (desarrollo rural, Cork 1996).
\item \autor{Fischler}: \concepto{pago único} desacoplado (2000-2002), \concepto{condicionalidad} obligatoria, \concepto{modulación}.
\item 2005: FEAGA/FEADER; 2007: OCM única; 2008: chequeo médico.
\end{itemize}
\subsub{d) La reforma de 2013}
\begin{itemize}
\item Pago básico con convergencia; \concepto{pago verde} (30\,\%; según el TCE, cambió prácticas en $\approx$5\,\% de la superficie).
\item Jóvenes, agricultor activo; fin de cuotas (leche 2015, azúcar 2017).
\end{itemize}

\sub{II.2. La PAC 2023-2027 y su simplificación\label{rep:II.2}}
\subsub{a) El nuevo modelo de aplicación}
\begin{itemize}
\item Reglamentos 2021/2115, 2116~y 2117; \concepto{plan estratégico} único por Estado (DAFO, metas, indicadores); 10~objetivos.
\item Pagos directos:
  \begin{itemize}
  \item ayuda básica (convergencia al 85\,\% en 2026); redistributivo $\geq$10\,\%; jóvenes $\geq$3\,\%;
  \item \concepto{ecorregímenes} $\geq$25\,\%; FEADER $\geq$35\,\% ambiental; reserva agrícola $\geq$450~M€.
  \end{itemize}
\end{itemize}
\subsub{b) Arquitectura verde y condicionalidad social}
\begin{itemize}
\item \concepto{Condicionalidad reforzada} (11~RLG + 9~BCAM) y \concepto{social} (obligatoria desde 2025).
\end{itemize}
\graficorepaso[0.9\linewidth]{arquitectura-verde}{Condicionalidad (obligatoria), ecorregímenes (anuales) y compromisos del FEADER (plurianuales)}
\subsub{c) La simplificación de 2024 y 2025}
\begin{itemize}
\item Reglamento 2024/1468: BCAM 6-8 suavizadas; $<$10~ha sin controles ($\approx$65\,\% de perceptores).
\item Diálogo estratégico (\autor{Strohschneider}, 2024); \emph{Visión} (19-2-2025): «incentivos antes que condiciones».
\item Reglamento 2025/2649: pequeños hasta 3.000 € sin condicionalidad; un control al año; ahorro 1.580~M€/año.
\end{itemize}
\subsub{d) España en la PAC}
\begin{itemize}
\item PEPAC 34.472~M€ (31.508~UE): directos 24.440, sectorial 1.712, rural 8.320; ecorregímenes 5.553 (23\,\%).
\item FEGA 2025: FEAGA 5.719 ($\approx$14\,\%), FEADER UE $\approx$1.680.
\end{itemize}

\sub{II.3. Problemas actuales y retos futuros\label{rep:II.3}}
\subsub{a) La distribución de las ayudas}
\begin{itemize}
\item 2023: el 46,5\,\% cobra $<$1.250 € (4,2\,\% del importe); el 25\,\% se lleva el 84,6\,\%; el 0,5\,\% ($>$100.000 €), el 15,9\,\%.
\item España: 20,7\,\% $\to$ 75,9\,\%. Pagos $\neq$ renta.
\end{itemize}
\graficorepaso[0.9\linewidth]{distribucion-pagos}{\% de perceptores frente a \% del importe; Lorenz: Gini $\geq$0,75 (UE) y 0,70 (España)}
\subsub{b) La capitalización de las ayudas en la tierra}
\begin{itemize}
\item \emph{Idea}: renta ricardiana \vertema{3.A.2}.
\item \emph{Supuestos}: mercado de arrendamiento, oferta rígida, pago $s$ por ha al arrendatario; $P_T$ = valor actual de las rentas.
\item \emph{Desarrollo}: la demanda sube $s$; incidencia \vertema{4.B.8}:
\funcionclave{\frac{\Delta R}{s}=\frac{|\varepsilon_D|}{|\varepsilon_D|+\varepsilon_S},\quad \Delta P_T=\frac{\Delta R}{i}}
\item \emph{Implicaciones}: gana el propietario; barrera de entrada; trampa de las ganancias transitorias (\autor{Tullock}).
\item \emph{Evidencia}: \autor{Kirwan} (2009): 25~céntimos por dólar; \autor{Ciaian}-\autor{Kancs} (2012): 0,19; con la reforma de 2013, 72-76\,\% (nuevos) y 20-39\,\% (antiguos).
\end{itemize}
\graficorepaso[0.85\linewidth]{capitalizacion-tierra}{Oferta rígida: $\Delta R=s$; con pendiente, parte del pago queda al agricultor}
\subsub{c) Relevo, medio ambiente y clima}
\begin{itemize}
\item 8,8~M de explotaciones (2023), $-$40\,\% desde 2005; $<$40~años 11,9\,\%, $\geq$65 33,2\,\% (España 7,7~y 41,3\,\%).
\item Tierra arable 11.000 $\to$ 15.200 €/ha (2018-2024).
\item TCE (2021): 100.000~M€ «climáticos» sin bajar emisiones; retirada del reglamento de plaguicidas (2024) \vertema{3.B.31}.
\end{itemize}
\subsub{d) Comercio exterior y ampliación}
\begin{itemize}
\item Caja verde y fin de restituciones (Nairobi, 2015) \vertema{3.B.33}.
\item Ucrania: medidas autónomas 2022-2025 y ALCAP revisado (2025); ampliación ($\approx$96.500~M€ en 7~años, según la prensa).
\item Mercosur, provisional desde 1-5-2026 con salvaguardia agrícola \vertema{3.B.35}.
\end{itemize}
\subsub{e) La PAC después de 2027}
\begin{itemize}
\item Propuesta (16-7-2025): sin pilares ni FEAGA/FEADER; Fondo de Planes Nacionales y Regionales (865.000~M€).
  \begin{itemize}
  \item 300.000~reservados (293.700~renta + 6.300~crisis); degresividad 25/50/75\,\% y tope 100.000.
  \end{itemize}
\item PE: 385.100~y 2~pilares; Consejo: 261.000-265.000; riesgo: renacionalización \vertema{3.B.37}.
\end{itemize}

% ==========================================================================
\bloque[8]{III. La política pesquera común\label{rep:III}}
\begin{ideaclave}
La pesca explota un recurso común: el libre acceso disipa la renta y agota la población. La PPC pasó de repartir el acceso (1983) a buscar el RMS (2013): la presión ha bajado en el Atlántico, pero el objetivo de 2020 no se cumplió y el Mediterráneo va tarde.
\end{ideaclave}

\sub{III.1. Justificación, naturaleza y objetivos\label{rep:III.1}}
\subsub{a) Relevancia y justificación}
\begin{itemize}
\item UE: 3,3~M de t (2025; 5,9~en 2001); $\approx$68.000~buques; autoabastecimiento 38\,\%.
\item España: 31\,\% del valor, 21\,\% del peso, 24\,\% del arqueo (8.324~buques).
\item Recurso común (rival, no excluible): tragedia de los comunes (\autor{Hardin}, 1968) \vertema{3.A.23}; 200~millas (1977).
\end{itemize}
\subsub{b) El modelo de \autor{Gordon} y \autor{Schaefer} (1954)}
\begin{itemize}
\item \emph{Supuestos}: $G(X)=\gamma X(1-X/\bar X)$; $H=\theta EX$; $p$ y $c$ dados; estados estacionarios.
\item \emph{Desarrollo}:
  \begin{itemize}
  \item $X(E)=\bar X(1-\theta E/\gamma)$; $Y(E)=\theta\bar XE(1-\theta E/\gamma)$;
  \item RMS: $E=\gamma/2\theta$, $X=\bar X/2$; libre acceso ($pY=cE$): $X_{\mathit{LA}}=c/p\theta$;
  \item MRE ($pY'(E)=c$):
\funcionclave{E_{\mathit{MRE}}=\tfrac12E_{\mathit{LA}}<E_{\mathit{RMS}},\quad X_{\mathit{MRE}}>X_{\mathit{RMS}}}
  \end{itemize}
\end{itemize}
\graficorepaso[0.95\linewidth]{gordon-schaefer}{Crecimiento y captura (izq.); ingreso y coste (der.): MRE $<$ RMS $<$ LA en esfuerzo}
\begin{itemize}
\item \emph{Implicaciones}: libre acceso doblemente ineficiente; RMS $\neq$ óptimo económico; subvenciones que bajan $c$ $\Rightarrow$ sobrepesca \vertema{3.B.33}.
\item \emph{Extensiones}: descuento (\autor{Clark}, 1990); TAC (carrera), \concepto{CIT} (\autor{Coase}), \autor{Ostrom}; poblaciones compartidas (\autor{Levhari}-\autor{Mirman}, 1980).
\item \emph{Evidencia}: FAO: 37,6\,\% no sostenibles (2023); CIT reducen a la mitad el colapso (\autor{Costello} \emph{et al.}, 2008).
\end{itemize}
\subsub{c) Competencias, marco jurídico y decisiones}
\begin{itemize}
\item Conservación exclusiva (art.~3.1.d); resto, compartido; el Consejo fija las posibilidades de pesca (art.~43.3).
\item Objetivos: Reglamento 1380/2013 (art.~2).
\end{itemize}

\sub{III.2. Análisis histórico: la PPC de 1983\label{rep:III.2}}
\begin{itemize}
\item 1970: igualdad de acceso y OCM; 200~millas (1977); CNUDM (1982).
\item Reglamento 170/83, 4~ejes:
  \begin{itemize}
  \item \emph{recursos}: \concepto{TAC} y cuotas con \concepto{estabilidad relativa}, 12~millas, medidas técnicas;
  \item \emph{estructuras}: IFOP $\to$ FEP $\to$ FEMP $\to$ FEMPA;
  \item \emph{mercados}: OCM, OP; \emph{exterior}: acuerdos, OROP.
  \end{itemize}
\item España (1986) duplica la capacidad; integración plena 1-1-1996.
\item Problemas: TAC políticos, sobrecapacidad, descartes.
\end{itemize}

\sub{III.3. Desarrollo y reformas: 2002 y 2013\label{rep:III.3}}
\begin{itemize}
\item \textbf{2002}: planes plurianuales, fin de ayudas a la construcción, consejos consultivos, INDNR (2008).
\item \textbf{2013} (Reg. 1380/2013): \concepto{RMS} en 2015-2020; \concepto{obligación de desembarque} (2015-2019).
  \begin{itemize}
  \item Regionalización; 11~consejos consultivos; ajuste de capacidad; etiquetado (Reg. 1379/2013).
  \end{itemize}
\end{itemize}

\sub{III.4. Situación actual\label{rep:III.4}}
\begin{itemize}
\item \emph{Atlántico NE}: $F/F_{\mathit{RMS}}$ 1,51 (2003) $\to$ 0,70 (2024); 81\,\% de las posibilidades de 2026 al RMS.
  \begin{itemize}
  \item Pero fuera de límites seguros 41 $\to$ 47\,\% (2021-2024).
  \item Solo el 29\,\% cumple el art.~2.2; caballa, 71\,\% por encima del ICES.
  \end{itemize}
\item \emph{Mediterráneo}: sobrepescadas 83\,\% (2007) $\to$ 47\,\% (2023); arrastre, 143~días en 2026 (la Comisión proponía 9,7).
\end{itemize}
\graficorepaso[0.85\linewidth]{estado-poblaciones}{$F/F_{\mathit{RMS}}$: el Atlántico, por debajo de 1~desde 2021; el Mediterráneo, en 2023}
\begin{itemize}
\item \emph{Exterior}: 11~protocolos; 5~OROP atuneras y 13~no atuneras; Brexit: 25\,\% de cuota (hasta 30-6-2026); acceso recíproco hasta 2038.
\item FEMPA: 6.108~M€; España 1.120 (21\,\%, 1.ª).
\item Control: Reg. 2023/2842, CATCH (2026), cámaras (2028).
\item Evaluación (abril de 2026): menos presión, RMS no alcanzado en 2020, el desembarque no cambió conductas, más importaciones.
\item Pacto por los Océanos (2025), Visión 2040 (2026), Ley del Océano (2027).
\end{itemize}

% ==========================================================================
\bloque[1]{Conclusión}
\begin{itemize}
\item \textbf{Recapitulación}: la PAC separó instrumentos (precio, pago, condiciones); la PPC pasó de repartir a conservar con el RMS.
\item \textbf{Relevancia}: casi 1/3 del presupuesto; España, gran beneficiaria.
\item \textbf{Extensiones}: \vertema{3.B.37} \vertema{3.B.35} \vertema{3.B.33} \vertema{3.B.31} \vertema{4.A.4}.
\item \textbf{Opinión}: pagar por bienes públicos, riesgo y relevo, no por hectárea; en el mar, RMS con derechos bien definidos.
\item \textbf{Idea final}: la mona sigue siendo mona si solo cambia el traje: hay que pagar por objetivos, no por tamaño.
\end{itemize}

% ==========================================================================
\anexorepaso
\bloque{Datos clave}
\begin{itemize}
\item PAC 2021-2027: \dato{386.600 M€}{FEAGA 291.100; FEADER 95.500}; \dato{0,31\,\% del PIB}{2024}.
\item Pagos directos: \dato{25\,\% $\to$ 84,6\,\%}{2023}.
\item PSE UE: \dato{38,5\,\% $\to$ 18,2\,\%}{OCDE, 1986-88 / 2025}.
\item Pesca: \dato{3,3 M de t}{2025}; $F/F_{\mathit{RMS}}$ Atlántico \dato{0,70}{2024}.
\end{itemize}
\bloque{Autores y fechas}
\hito{1954}{\autor{Gordon}, \autor{Schaefer}: pesca}
\hito{1958}{Stresa; \autor{Cochrane}: rueda del progreso}
\hito{1962}{OCM y FEOGA}
\hito{1965}{\autor{Olson}: acción colectiva}
\hito{1983}{Primera PPC}
\hito{1992}{Reforma \autor{MacSharry}}
\hito{2003}{\autor{Fischler}: pago único}
\hito{2013}{PAC verde; PPC del RMS}
\bloque{Preguntas del tribunal}
\begin{enumerate}
\item \textbf{¿Por qué el pago desacoplado no es neutral?} Riqueza, crédito, expectativas, salida y capitalización.
\item \textbf{¿Quién se queda con el pago por hectárea?} Según las elasticidades; con oferta rígida, el propietario.
\item \textbf{¿RMS o MRE?} El MRE deja más peces con menos esfuerzo; el RMS es biológico.
\item \textbf{¿TAC o CIT?} El TAC no frena la carrera; las CIT dan un precio a la renta, pero concentran.
\item \textbf{¿Fondo único tras 2027?} Flexibilidad y resultados frente a renacionalización.
\end{enumerate}
\begin{trampa}
\begin{itemize}
\item Conservación de los recursos marinos: \emph{exclusiva}.
\item \autor{MacSharry}: caja \emph{azul}; desacoplados: caja \emph{verde}.
\item Ecorregímenes $\geq$\emph{25}\,\%; redistributivo $\geq$\emph{10}\,\%.
\item $E_{\mathit{MRE}}=E_{\mathit{LA}}/2$; biomasa del RMS, $\bar X/2$.
\end{itemize}
\end{trampa}

\end{columnas}
\end{document}
